\subsection*{D7. Offline selection, stopping and labor accounting}
\label{supp:d7-labor}
\label{sec:labor-history}

These comparisons use annotator histories for re-annotation. Target labor is recorded annotation person-time; project labor also includes history acquisition. C8--C9 specify the policies and accounting.

At team size 5, quality-and-time selection gives mean $Q=93.926$ versus 93.708 for quality-only selection across 126 contexts from 46 images and 9 buildings. The paired difference is $+0.218$ (95\% building-bootstrap CI $[-0.538,0.869]$), with a geometric labor ratio of 0.813 ($[0.760,0.882]$). Independently selected 13-person teams require more labor and have lower mean $Q$ than 5-person teams under both policies on common support (C13).

Source-history stopping raises mean $Q$ from 85.189 to 86.054 relative to fixed-7 teams, improving 13 of 21 images in 5 room groups. On the separate timing panel of 11 images in 3 groups, target labor rises from 1,307.9 to 1,660.3 person-seconds. Including history acquisition gives a project-labor ratio of 2.115. The stopping rule improves observed quality and requires more work.
